\begin{table}[t]
\centering
\caption{Standardized Effect Sizes}
\label{tab:sde}
\begin{tabular}{lcccccc}
\toprule
Outcome & $\hat{\beta}$ & SE & SD($Y$) & SDE & SE(SDE) & Classification \\
\midrule
\multicolumn{7}{l}{\textit{Panel A: Pooled}} \\
Wage gap ratio & 2.20 & 0.22 & 5.60 & 0.393 & 0.040 & Large positive \\
Non-regular wage & -3.54 & 0.06 & 10.15 & -0.349 & 0.006 & Large negative \\
Regular wage & -12.91 & 0.65 & 42.27 & -0.305 & 0.015 & Large negative \\
\midrule
\multicolumn{7}{l}{\textit{Panel B: Heterogeneous (by sex)}} \\
Gap ratio (male) & 0.55 & 0.08 & 6.26 & 0.089 & 0.013 & Moderate positive \\
Gap ratio (female) & 3.02 & 0.06 & 4.28 & 0.705 & 0.014 & Large positive \\
\bottomrule
\end{tabular}
\begin{tablenotes}
\item \textit{Notes:} \textbf{Country:} Japan. \textbf{Research question:} Does anti-discrimination legislation (the 2020 Equal Pay for Equal Work Act) narrow the wage gap between regular and non-regular workers in a dual labor market? \textbf{Policy mechanism:} The revised Part-Time/Fixed-Term Employment Act prohibits unreasonable differences in wages, bonuses, and benefits between regular and non-regular employees performing similar work, enforced through administrative guidance and litigation rather than statutory penalties. \textbf{Outcome definition:} Wage gap ratio defined as (non-regular monthly scheduled wage / regular monthly scheduled wage) $\times$ 100, where higher values indicate a narrower gap. \textbf{Treatment:} Binary, staggered by firm size: large firms ($\geq$300 employees) treated April 2020, SMEs ($<$300) treated April 2021. \textbf{Data:} MHLW Basic Survey on Wage Structure, 2014--2024, firm-size $\times$ sex cells, 33 observations (total sex panel). \textbf{Method:} Callaway--Sant'Anna (2021) staggered DiD with not-yet-treated control group, clustered by firm size. \textbf{Sample:} All general workers (excluding part-time) in private establishments with 10+ employees; 3 firm-size categories $\times$ 11 annual surveys. SDE $= \hat{\beta} / \text{SD}(Y)$ where SD($Y$) is the pre-treatment standard deviation. Classification refers to magnitude, not statistical significance: Large ($|$SDE$| > 0.15$), Moderate ($0.05$--$0.15$), Small ($0.005$--$0.05$), Null ($< 0.005$).
\end{tablenotes}
\end{table}
